\ifdefined\gkver\else\def\gkver{student}\fi
\documentclass[\gkver]{gaokaozhenti}
\begin{document}

\chapter{2015年上海}

\begin{enumerate}
\item
%% number: 1
%% paperName: 2015年上海高考第 1 题 2 分
%% typeId: 1
%% type: 单选题
%% chapter: 光学
%% point: 光的电磁说
%% method: 无
%% score: 2
%% degree: 843
%% duplicateId: 0
%% body:
X 射线\xzanswer{B}

\fourchoices[answer=B]
{不是电磁波}
{具有反射和折射的特性}
{只能在介质中传播}
{不能发生干涉和衍射}

%% answer: B
%% memo:
\memoanswer{
X 射线是电磁波，具有波的性质，可以发生反射、折射、干涉和衍射，可以在真空中传播，故 B 正确 ACD 错误。
}

\item
%% number: 2
%% paperName: 2015年上海高考第 2 题 2 分
%% typeId: 1
%% type: 单选题
%% chapter: 机械波
%% point: 干涉衍射
%% method: 无
%% score: 2
%% degree: 844
%% duplicateId: 0
%% body:
如图，P 为桥墩，A 为靠近桥墩浮在水面的叶片，波源 S 连续振动，形成水波，此时叶片 A 静止不动。为使水波能带动叶片振动，可用的方法是\xzanswer{B}

\onepicture[h=3.0cm]{\includesvg{figs/02}}

\fourchoices[answer=B]
{提高波源频率}
{降低波源频率}
{增加波源距桥墩的距离}
{减小波源距桥墩的距离}

%% answer: B
%% memo:
\memoanswer{
A 静止不动的原因是衍射现象不明显，要想衍射现象变得明显，需要增大波长，由 $v=\lambda f$ 可得，介质不变，波速不变，要想增大波长，只能降低频率，故 B 正确。
}

\item
%% number: 3
%% paperName: 2015年上海高考第 3 题 2 分
%% typeId: 1
%% type: 单选题
%% chapter: 运动和力
%% point: 牛顿第二定律
%% method: 无
%% score: 2
%% degree: 758
%% duplicateId: 0
%% body:
如图，鸟沿虚线斜向上加速飞行，空气对其作用力可能是\xzanswer{B}

\onepicture[h=3.0cm]{\includesvg{figs/03}}

\fourchoices[answer=B]
{$F_{1}$}
{$F_{2}$}
{$F_{3}$}
{$F_{4}$}

%% answer: B
%% memo:
\memoanswer{
空气对鸟的作用力，一方面抵消重力，另一方面让鸟沿虚线加速飞行，所以空气对鸟的作用力与鸟的重力的合力沿虚线斜向上，空气对鸟的作用力可能是 $F_{2}$，故 B 正确。

\onepicture[h=2.8cm]{figs/03_an.png}
}

\item
%% number: 4
%% paperName: 2015年上海高考第 4 题 2 分
%% typeId: 1
%% type: 单选题
%% chapter: 气体性质
%% point: 物体的内能
%% method: 无
%% score: 2
%% degree: 990
%% duplicateId: 0
%% body:
一定质量的理想气体在升温过程中\xzanswer{C}

\fourchoices[answer=C]
{分子平均势能减小}
{每个分子速率都增大}
{分子平均动能增大}
{分子间作用力先增大后减小}

%% answer: C
%% memo:
\memoanswer{
理想气体不考虑分子势能，不考虑分子间的作用力，故 AD 错误；温度升高，分子的平均动能增大，分子的平均速率增大，但不是每个分子的速率都增大，故 B 错误 C 正确。
}

\item
%% number: 5
%% paperName: 2015年上海高考第 5 题 2 分
%% typeId: 1
%% type: 单选题
%% chapter: 原子和原子核
%% point: 核裂变
%% method: 无
%% score: 2
%% degree: 918
%% duplicateId: 0
%% body:
铀核可以发生衰变和裂变，铀核的\xzanswer{C}

\fourchoices[answer=C]
{衰变和裂变都能自发发生}
{衰变和裂变都不能自发发生}
{衰变能自发发生而裂变不能自发发生}
{衰变不能自发发生而裂变能自发发生}

%% answer: C
%% memo:
\memoanswer{
铀核的衰变是可以自发产生，但裂变需要一定的条件，故 C 正确。
}

\item
%% number: 6
%% paperName: 2015年上海高考第 6 题 2 分
%% typeId: 1
%% type: 单选题
%% chapter: 原子和原子核
%% point: 天然放射性
%% method: 无
%% score: 2
%% degree: 934
%% duplicateId: 0
%% body:
$\ce{^{232}_{90}Th}$ 经过一系列 $\alpha$ 衰变和 $\beta$ 衰变后变成 $\ce{^{208}_{82}Pb}$，则 $\ce{^{208}_{82}Pb}$ 比 $\ce{^{232}_{90}Th}$ 少\xzanswer{A}

\fourchoices[answer=A]
{16 个中子，8 个质子}
{8 个中子，16 个质子}
{24 个中子，8 个质子}
{8 个中子，24 个质子}

%% answer: A
%% memo:
\memoanswer{
由反应前后质量数和电荷数守恒，则 $\ce{^{208}_{82}Pb}$ 比 $\ce{^{232}_{90}Th}$ 少的质子数为 $90-82=8$，少的中子数为 $(232-90)-(208-82)=16$，故 A 正确。
}

\item
%% number: 7
%% paperName: 2015年上海高考第 7 题 2 分
%% typeId: 1
%% type: 单选题
%% chapter: 原子和原子核
%% point: 原子核式结构
%% method: 无
%% score: 2
%% degree: 736
%% duplicateId: 0
%% body:
在 $\alpha$ 粒子散射实验中，电子对 $\alpha$ 粒子运动的影响可以忽略。这是因为与 $\alpha$ 粒子相比，电子的\xzanswer{D}

\fourchoices[answer=D]
{电量太小}
{速度太小}
{体积太小}
{质量太小}

%% answer: D
%% memo:
\memoanswer{
电子的质量比 $\alpha$ 粒子小得多，约为 $\alpha$ 粒子质量的 $\frac{1}{7400}$，所以对 $\alpha$ 粒子运动的影响可以忽略，故 D 正确。
}

\item
%% number: 8
%% paperName: 2015年上海高考第 8 题 2 分
%% typeId: 1
%% type: 单选题
%% chapter: 电场
%% point: 电场线
%% method: 无
%% score: 2
%% degree: 651
%% duplicateId: 0
%% body:
两个正、负点电荷周围电场线分布如图所示。P、Q 为电场中两点，则\xzanswer{D}

\onepicture[h=3.4cm]{\includesvg{figs/08}}

\fourchoices[answer=D]
{正电荷由 P 静止释放能运动到 Q}
{正电荷在 P 的加速度小于在 Q 的加速度}
{负电荷在 P 的电势能高于在 Q 的电势能}
{负电荷从 P 移动到 Q，其间必有一点电势能为零}

%% answer: D
%% memo:
\memoanswer{
电荷在匀强电场中静止释放才会沿着电场线运动，正电荷由 P 点释放，将沿 P 点的切线方向运动，不能到达 Q 点，故 A 错误；P 点的电场线比 Q 点密，故 P 点的场强较大，正电荷在 P 点的加速度大于在 Q 点的加速度，故 B 错误；负电荷在电势高的位置电势能小，P 点电势高于 Q 点电势，所以负电荷在 P 点的电势能小于在 Q 点的电势能，故 C 错误；P 点电势大于零，Q 点电势小于零，P、Q 之间一定有一点电势为零，负电荷在该点的电势能也为零，故 D 正确。
}

\item
%% number: 9
%% paperName: 2015年上海高考第 9 题 3 分
%% typeId: 1
%% type: 单选题
%% chapter: 气体性质
%% point: 液柱移动定性分析
%% method: 无
%% score: 3
%% degree: 709
%% duplicateId: 0
%% body:
如图，长为 $h$ 的水银柱将上端封闭的玻璃管内气体分隔成两部分，A 处管内外水银面相平。将玻璃管缓慢向上提升 $H$ 高度（管下端未离开水银面），上下两部分气体的压强变化分别为 $\Delta p_{1}$ 和 $\Delta p_{2}$，体积变化分别为 $\Delta V_{1}$ 和 $\Delta V_{2}$。已知水银密度为 $\rho$，玻璃管截面积为 $S$，则\xzanswer{A}

\onepicture[h=3.6cm]{\includesvg{figs/09}}

\fourchoices[answer=A]
{$\Delta p_{2}$ 一定等于 $\Delta p_{1}$}
{$\Delta V_{2}$ 一定等于 $\Delta V_{1}$}
{$\Delta p_{2}$ 与 $\Delta p_{1}$ 之差为 $\rho gh$}
{$\Delta V_{2}$ 与 $\Delta V_{1}$ 之和为 $HS$}

%% answer: A
%% memo:
\memoanswer{
上下两部分气体的体积关系不确定，两部分气体均发生等温变化，则 $\Delta V_{1}$ 和 $\Delta V_{2}$ 不一定相等，提升玻璃管，再次平衡后，玻璃管中的水银面高于水银槽中的水银面，但高度之差小于 $H$，故 $\Delta V_{1}$ 和 $\Delta V_{2}$ 之和小于 $HS$，故 BD 错误；提升之前：$p_{2}-p_{1}=\rho gh$，提升之后 $p_{2}'-p_{1}'=\rho gh$，联立可得 $\Delta p_{1}=\Delta p_{2}$，故 A 正确 C 错误。
}

\item
%% number: 10
%% paperName: 2015年上海高考第 10 题 3 分
%% typeId: 1
%% type: 单选题
%% chapter: 光学
%% point: 波粒二象性
%% method: 无
%% score: 3
%% degree: 578
%% duplicateId: 0
%% body:
用很弱的光做单缝衍射实验，改变曝光时间，在胶片上出现的图像如图所示，该实验表明\xzanswer{C}

\onepicture[h=2.2cm]{\includesvg{figs/10}}

\fourchoices[answer=C]
{光的本质是波}
{光的本质是粒子}
{光的能量在胶片上分布不均匀}
{光到达胶片上不同位置的概率相同}

%% answer: C
%% memo:
\memoanswer{
光的本质是具有波粒二象性，少量光子具有粒子性，大量光子具有波动性，故 AB 错误；由图像可以看出，光的能量在胶片上分布不均匀，光到达胶片上不同位置的概率不同，故 C 正确 D 错误。
}

\item
%% number: 11
%% paperName: 2015年上海高考第 11 题 3 分
%% typeId: 1
%% type: 单选题
%% chapter: 光学
%% point: 光电效应
%% method: 无
%% score: 3
%% degree: 441
%% duplicateId: 0
%% body:
某光源发出的光由不同波长的光组成，不同波长的光的强度如图所示。表中给出了一些材料的极限波长，用该光源发出的光照射表中材料\xzanswer{D}

\onepicture[h=2.6cm]{\includesvg{figs/11}}

\begin{center}
\begin{tabular}{|c|c|c|c|}
\hline
材料 & 钠 & 铜 & 铂 \\
\hline
极限波长（$\Unm$） & 541 & 268 & 196 \\
\hline
\end{tabular}
\end{center}

\fourchoices[answer=D]
{仅钠能产生光电子}
{仅钠、铜能产生光电子}
{仅铜、铂能产生光电子}
{都能产生光电子}

%% answer: D
%% memo:
\memoanswer{
波长越长，频率越小，故波长比材料极限波长小的光照射该材料时可以发生光电效应，结合图象和表格，用该光源照射表中材料都能产生光电子，故 D 正确。
}

\item
%% number: 12
%% paperName: 2015年上海高考第 12 题 3 分
%% typeId: 1
%% type: 单选题
%% chapter: 电路
%% point: 电流强度
%% method: 无
%% score: 3
%% degree: 721
%% duplicateId: 0
%% body:
重离子肿瘤治疗装置中的回旋加速器可发射 $+5$ 价重离子束，其电流强度为 $1.2\times10^{-5}\UA$，则在 $1\Us$ 内发射的重离子个数为（$e=1.6\times10^{-19}\UC$）\xzanswer{B}

\fourchoices[answer=B]
{$3.0\times10^{12}$}
{$1.5\times10^{13}$}
{$7.5\times10^{13}$}
{$3.75\times10^{14}$}

%% answer: B
%% memo:
\memoanswer{
$1\Us$ 内发射的重离子个数为 $n=\frac{It}{5e}=\frac{1.2\times10^{-5}\times1}{5\times1.6\times10^{-19}}=1.5\times10^{13}$，故 B 正确。
}

\item
%% number: 13
%% paperName: 2015年上海高考第 13 题 3 分
%% typeId: 1
%% type: 单选题
%% chapter: 电路
%% point: 逻辑电路
%% method: 无
%% score: 3
%% degree: 821
%% duplicateId: 0
%% body:
监控系统控制电路如图所示，电键 S 闭合时，系统白天和晚上都工作；电键 S 断开时，系统仅晚上工作。在电路中虚框处分别接入光敏电阻（受光照时阻值减小）和定值电阻，则电路中\xzanswer{D}

\onepicture[h=3.6cm]{\includesvg{figs/13}}

\fourchoices[answer=D]
{C 是“与门”，A 是光敏电阻}
{C 是“与门”，B 是光敏电阻}
{C 是“或门”，A 是光敏电阻}
{C 是“或门”，B 是光敏电阻}

%% answer: D
%% memo:
\memoanswer{
要想使监控系统白天和晚上都工作，C 是“或门”，S 断开时，系统仅晚上工作，说明 B 是光敏电阻，故选项 D 正确。
}

\item
%% number: 14
%% paperName: 2015年上海高考第 14 题 3 分
%% typeId: 1
%% type: 单选题
%% chapter: 力物体的平衡
%% point: 多力平衡
%% method: 无
%% score: 3
%% degree: 655
%% duplicateId: 0
%% body:
如图，一质量为 $m$ 的正方体物块置于风洞内的水平面上，一面与风速垂直，当风速为 $v_{0}$ 时刚好能推动该物块。已知风对物块的推力 $F\propto Sv^{2}$，其中 $v$ 为风速、$S$ 为物块迎风面积。当风速变为 $2v_{0}$ 时，刚好能推动用同一材料做成的另一正方体物块，则该物块的质量为\xzanswer{D}

\onepicture[h=2.4cm]{\includesvg{figs/14}}

\fourchoices[answer=D]
{$4m$}
{$8m$}
{$32m$}
{$64m$}

%% answer: D
%% memo:
\memoanswer{
设正方体物块边长为 $L$，该材料的密度为 $\rho$，与地面的动摩擦因数为 $\mu$，由共点力平衡有 $\mu\rho L^{3}g=kL^{2}v^{2}$，即 $kv^{2}=\mu\rho Lg$，速度变为原来的 2 倍，则边长应变为原来的 4 倍，由 $m=\rho L^{3}$ 可知质量应为原来的 64 倍，故 D 正确。
}

\item
%% number: 15
%% paperName: 2015年上海高考第 15 题 3 分
%% typeId: 1
%% type: 单选题
%% chapter: 机械波
%% point: 波的图象
%% method: 无
%% score: 3
%% degree: 728
%% duplicateId: 0
%% body:
一简谐横波沿水平绳向右传播，波速为 $v$，周期为 $T$，振幅为 $A$。绳上两质点 M、N 的平衡位置相距 $\frac{3}{4}$ 波长，N 位于 M 右方。设向上为正，在 $t=0$ 时 M 位移为 $+\frac{A}{2}$，且向上运动；经时间 $t$（$t<T$），M 位移仍为 $+\frac{A}{2}$，但向下运动，则\xzanswer{C}

\fourchoices[answer=C]
{在 $t$ 时刻，N 恰好在波谷位置}
{在 $t$ 时刻，N 位移为负，速度向上}
{在 $t$ 时刻，N 位移为负，速度向下}
{在 $2t$ 时刻，N 位移为 $-\frac{A}{2}$，速度向下}

%% answer: C
%% memo:
\memoanswer{
由题可得 $t'=0$ 时、$t'=t$ 时绳上的大致波形如图，$t'=0$ 时波形如实线所示，$t'=t$ 时刻的波形如虚线所示，可知在 $t$ 时刻，N 的位移为负，速度向下，故 AB 错误 C 正确；在 $2t$ 时刻，N 的速度方向向上，故 D 错误。

\onepicture[h=2.4cm]{figs/15_an.jpg}
}

\item
%% number: 16
%% paperName: 2015年上海高考第 16 题 3 分
%% typeId: 1
%% type: 单选题
%% chapter: 曲线运动
%% point: 平抛运动
%% method: 无
%% score: 3
%% degree: 322
%% duplicateId: 0
%% body:
如图，战机在斜坡上方进行投弹演练。战机水平匀速飞行，每隔相等时间释放一颗炸弹，第一颗落在 a 点，第二颗落在 b 点。斜坡上 c、d 两点与 a、b 共线，且 $ab=bc=cd$，不计空气阻力。第三颗炸弹将落在\xzanswer{A}

\onepicture[h=2.6cm]{\includesvg{figs/16}}

\fourchoices[answer=A]
{bc 之间}
{c 点}
{cd 之间}
{d 点}

%% answer: A
%% memo:
\memoanswer{
所有炸弹在水平方向做匀速直线运动，速度大小是相同的，相同时间间隔第三颗炸弹落在斜面上时，竖直方向下落的距离大，故第三颗炸弹将落在 b、c 之间，故 A 正确。
}

\item
%% number: 17
%% paperName: 2015年上海高考第 17 题 4 分
%% typeId: 2
%% type: 多选题
%% chapter: 机械振动
%% point: 振动图象
%% method: 无
%% score: 4
%% degree: 532
%% duplicateId: 0
%% body:
质点运动的位移 $x$ 与时间 $t$ 的关系如图所示，其中做机械振动的是\xzanswer{ABC}

\fourchoices[answer=ABC, ispicture=true, h=2.2cm]
{\includesvg[inkscapearea=page]{figs/17a}}
{\includesvg[inkscapearea=page]{figs/17b}}
{\includesvg[inkscapearea=page]{figs/17c}}
{\includesvg[inkscapearea=page]{figs/17d}}

%% answer: ABC
%% memo:
\memoanswer{
在平衡位置附近的往复运动是机械振动，故 ABC 正确。
}

\item
%% number: 18
%% paperName: 2015年上海高考第 18 题 4 分
%% typeId: 2
%% type: 多选题
%% chapter: 曲线运动
%% point: 竖直平面圆周运动
%% method: 对称法
%% score: 4
%% degree: 209
%% duplicateId: 0
%% body:
如图，质量为 $m$ 的小球用轻绳悬挂在 O 点，在水平恒力 $F=mg\tan\theta$ 作用下，小球从静止开始由 A 经 B 向 C 运动。则小球\xzanswer{ACD}

\onepicture[h=4.4cm]{\includesvg{figs/18}}

\fourchoices[answer=ACD]
{先加速后减速}
{在 B 点加速度为零}
{在 C 点速度为零}
{在 C 点加速度为 $g\tan\theta$}

%% answer: ACD
%% memo:
\memoanswer{
对小球受力分析，在垂直于绳的方向有 $mg\tan\theta\cos\alpha-mg\sin\alpha=ma$，$\alpha$ 为绳与竖直方向的夹角，可得由 A 到 B：$a>0$，由 B 到 C：$a<0$，所以小球由 A 到 B 做加速运动，B 到 C 做减速运动，故 A 正确；在 B 点垂直绳方向的合力为零，但沿绳方向的合力不为零，小球具有向心加速度，故 B 错误；C 点是与 A 点对称的位置，速度为零，故 C 正确；在 C 点的加速度大小为 $\left|\frac{mg\tan\theta\cos2\theta-mg\sin2\theta}{m}\right|=g\tan\theta$，故 D 正确。
}

\item
%% number: 19
%% paperName: 2015年上海高考第 19 题 4 分
%% typeId: 2
%% type: 多选题
%% chapter: 功和能
%% point: 子弹打木块模型
%% method: 图象法
%% score: 4
%% degree: 342
%% duplicateId: 0
%% body:
一颗子弹以水平速度 $v_{0}$ 穿透一块在光滑水平面上迎面滑来的木块后，二者运动方向均不变。设子弹与木块间相互作用力恒定，木块最后速度为 $v$，则\xzanswer{AC}

\fourchoices[answer=AC]
{$v_{0}$ 越大，$v$ 越大}
{$v_{0}$ 越小，$v$ 越大}
{子弹质量越大，$v$ 越大}
{木块质量越小，$v$ 越大}

%% answer: AC
%% memo:
\memoanswer{
子弹穿过木块的过程中，两者间的相互作用力使木块和子弹都做减速运动，$v_{0}$ 越大，子弹穿过的时间越短。木块减速时间越短，由 $v=v_{\text{块}}-at$ 知 $v$ 越大，故 A 正确，B 错误；子弹质量越大，子弹的加速度越小，在木块中减速运动时间越短，同理知 $v$ 越大，故 C 正确；木块质量越小，加速度越大，同理知 $v$ 越小，故 D 错误。
}

\item
%% number: 20
%% paperName: 2015年上海高考第 20 题 4 分
%% typeId: 2
%% type: 多选题
%% chapter: 电磁感应
%% point: 电磁感应能量分析
%% method: 无
%% score: 4
%% degree: 419
%% duplicateId: 0
%% body:
如图，光滑平行金属导轨固定在水平面上，左端由导线相连，导体棒垂直静置于导轨上构成回路。在外力 $F$ 作用下，回路上方的条形磁铁竖直向上做匀速运动。在匀速运动过程中外力 $F$ 做功 $W_{F}$，磁场力对导体棒做功 $W_{1}$，磁铁克服磁场力做功 $W_{2}$，重力对磁铁做功 $W_{G}$，回路中产生的焦耳热为 $Q$，导体棒获得的动能为 $E_{k}$。则\xzanswer{BCD}

\onepicture[h=3.6cm]{\includesvg{figs/20}}

\fourchoices[answer=BCD]
{$W_{1}=0$}
{$W_{2}-W_{1}=Q$}
{$W_{1}=E_{k}$}
{$W_{F}+W_{G}=Q+E_{k}$}

%% answer: BCD
%% memo:
\memoanswer{
由题知，磁场力对导体棒做的功等于导体棒获得的动能，即 $W_{1}=E_{k}$，故 C 正确 A 错误；又 $W_{2}=Q+E_{k}$，故 $W_{2}-W_{1}=Q$，故 B 正确；由功能关系，对磁铁 $W_{F}-W_{2}+W_{G}=0$，联立可得，故 D 正确。
}

\item
%% number: 21
%% paperName: 2015年上海高考第 21 题 4 分
%% typeId: 3
%% type: 填空题
%% chapter: 电场
%% point: 电势能电势差电场力做功
%% method: 无
%% score: 4
%% degree: 850
%% duplicateId: 0
%% body:
静电场是\tkanswer[2.0cm]{静止电荷}周围空间存在的一种物质；通常用\tkanswer[1.6cm]{电势}来描述电场的能的性质。
%% answer: 静止电荷，电势
%% memo:
\memoanswer{
静电场是一种客观物质，是静止电荷周围存在的一种特殊物质，通常用电势来描述电场的能的性质。
}

\item
%% number: 22
%% paperName: 2015年上海高考第 22 题 4 分
%% typeId: 3
%% type: 填空题
%% chapter: 曲线运动
%% point: 双星问题
%% method: 无
%% score: 4
%% degree: 450
%% duplicateId: 0
%% body:
（本大题中第 22 题为分叉题，分 A、B 两类，考生可任选一类答题。若两类试题均做，一律按 A 类题计分。）

22A．两小孩在冰面上乘坐“碰碰车”相向运动。A 车总质量为 $50\Ukg$，以 $2\Ums$ 的速度向右运动；B 车总质量为 $70\Ukg$，以 $3\Ums$ 的速度向左运动。碰撞后，A 以 $1.5\Ums$ 的速度向左运动，则 B 的速度大小为\tkanswer[1.4cm]{0.5}$\Ums$，方向向\tkanswer[1.2cm]{左}（选填“左”或“右”）。

22B．两靠得较近的天体组成的系统称为双星，它们以两者连线上某点为圆心做匀速圆周运动，因而不至于由于引力作用而吸引在一起。设两天体的质量分别为 $m_{1}$ 和 $m_{2}$，则它们的轨道半径之比 $R_{m1}:R_{m2}=$\tkanswer[2.0cm]{$m_{2}:m_{1}$}；速度之比 $v_{m1}:v_{m2}=$\tkanswer[2.0cm]{$m_{2}:m_{1}$}。
%% answer: 0.5，左；$m_{2}:m_{1}$，$m_{2}:m_{1}$
%% memo:
\memoanswer{
22A．规定向右为正方向，由动量守恒定律 $m_{A}v_{A}+m_{B}v_{B}=m_{A}v_{A}'+m_{B}v_{B}'$，解得 $v_{B}'=-0.5\Ums$，所以 B 的速度大小为 $0.5\Ums$，方向向左。

22B．双星的角速度相等，由 $\frac{Gm_{1}m_{2}}{r^{2}}=m_{1}\omega^{2}R_{m_{1}}=m_{2}\omega^{2}R_{m_{2}}$ 知，轨道半径与质量成反比，所以 $R_{m_{1}}:R_{m_{2}}=m_{2}:m_{1}$，由 $v=\omega R$ 知，线速度与半径成正比，所以 $v_{m_{1}}:v_{m_{2}}=m_{2}:m_{1}$。
}

\item
%% number: 23
%% paperName: 2015年上海高考第 23 题 4 分
%% typeId: 3
%% type: 填空题
%% chapter: 功和能
%% point: 功率
%% method: 无
%% score: 4
%% degree: 450
%% duplicateId: 0
%% body:
如图，汽车在平直路面上匀速运动，用跨过光滑定滑轮的轻绳牵引轮船，汽车与滑轮间的绳保持水平。当牵引轮船的绳与水平方向成 $\theta$ 角时，轮船速度为 $v$，绳的拉力对船做功的功率为 $P$，此时绳对船的拉力为\tkanswer[1.8cm]{$\frac{P}{v\cos\theta}$}。若汽车还受到恒定阻力 $f$，则汽车发动机的输出功率为\tkanswer[2.6cm]{$P+fv\cos\theta$}。

\onepicture[h=3.0cm, align=right]{\includesvg{figs/23}}
%% answer: $\frac{P}{v\cos\theta}$，$P+fv\cos\theta$
%% memo:
\memoanswer{
船速为 $v$，车的速度等于绳子的速度，大小为 $v\cos\theta$，绳子对船的拉力为 $F=\frac{P}{v\cos\theta}$，汽车发电机的输出功率为 $(F+f)v\cos\theta=\left(\frac{P}{v\cos\theta}+f\right)v\cos\theta=P+fv\cos\theta$。
}

\item
%% number: 24
%% paperName: 2015年上海高考第 24 题 4 分
%% typeId: 3
%% type: 填空题
%% chapter: 电磁感应
%% point: 电磁感应中的变加速
%% method: 无
%% score: 4
%% degree: 450
%% duplicateId: 0
%% body:
如图，一无限长通电直导线固定在光滑水平面上，金属环质量为 $0.02\Ukg$，在该平面上以 $v_{0}=2\Ums$、与导线成 $60^{\circ}$ 角的初速度运动，其最终的运动状态是\tkanswer[2.6cm]{匀速直线运动}，环中最多能产生\tkanswer[1.8cm]{$0.03\UJ$}的电能。

\onepicture[h=3.0cm, align=right]{\includesvg{figs/24}}
%% answer: 匀速直线运动，$0.03\UJ$
%% memo:
\memoanswer{
由无限长通电直导线周围的磁场分布可知，圆环在垂直导线的方向上受安培力，做减速运动，在沿导线方向受力平衡，做匀速直线运动，最终圆环沿导线方向做匀速直线运动，最多能产生的电能为 $W=\Delta E_{k}=\frac{1}{2}m(v_{0}\sin60^{\circ})^{2}=0.03\UJ$。
}

\item
%% number: 25
%% paperName: 2015年上海高考第 25 题 4 分
%% typeId: 3
%% type: 填空题
%% chapter: 磁场
%% point: 安培力
%% method: 对称法
%% score: 4
%% degree: 450
%% duplicateId: 0
%% body:
如图，两根通电长直导线 a、b 平行放置，a、b 中的电流强度分别为 $I$ 和 $2I$，此时 a 受到的磁场力为 $F$，若以该磁场力的方向为正，则 b 受到的磁场力为\tkanswer[1.2cm]{$-F$}。当在 a、b 的正中间再放置一根与 a、b 平行共面的通电长直导线 c 后，a 受到的磁场力大小变为 $2F$，则此时 b 受到的磁场力为\tkanswer[2.2cm]{$-3F$ 或 $5F$}。

\onepicture[h=3.8cm, align=right]{\includesvg{figs/25}}
%% answer: $-F$，$-3F$ 或 $5F$
%% memo:
\memoanswer{
a、b 之间的磁场力为作用力和反作用力，故 b 受到的磁场力大小也为 $F$，方向与 $F$ 方向相反，所以是 $-F$；若放 c 后 a 受到的磁场力为 $2F$，由题意，c 对 a 的磁场力为 $F$，那么，c 对 a 的作用力为 $-2F$，b 受到的磁场力为 $-3F$。若放 c 后 a 受到的磁场力为 $-2F$，由题意，c 对 a 的磁场力为 $-3F$，那么，c 对 b 的磁场力为 $6F$，b 受到的磁场力为 $5F$。
}

\item
%% number: 26
%% paperName: 2015年上海高考第 26 题 3 分
%% typeId: 4
%% type: 实验题
%% chapter: 磁场
%% point: 用DIS研究通电螺线管的磁感应强度
%% method: 无
%% score: 3
%% degree: 650
%% duplicateId: 0
%% body:
在“用 DIS 研究通电螺线管的磁感应强度”实验中
\begin{enumerate}
	\item
	在对螺线管通电\tkanswer[1.2cm]{前}（选填“前”或“后”）必须对磁传感器进行调零。

	\item
	（单选题）实验时，将磁传感器探管前端插至通电螺线管轴线中点时，磁传感器读数为 $5\UmT$。减小通电螺线管的电流后，将探管从螺线管的另一端插入，当探管前端再次到达螺线管轴线中点时，磁传感器的读数可能为

	\fourchoices[answer=D]
	{$5\UmT$}
	{$-5\UmT$}
	{$3\UmT$}
	{$-3\UmT$}
\end{enumerate}
%% answer: （1）前；（2）D
\jdanswer{
\begin{enumerate}
	\item
	前

	\item
	D

\end{enumerate}
}
%% memo:
\memoanswer{
（1）在通电前要对传感器进行调零；（2）减小电流，磁感应强度会变小，方向不变，由于探管从另一端插入，故 D 正确。
}

\item
%% number: 27
%% paperName: 2015年上海高考第 27 题 4 分
%% typeId: 4
%% type: 实验题
%% chapter: 电路
%% point: 多用表的使用
%% method: 无
%% score: 4
%% degree: 450
%% duplicateId: 0
%% body:
如图是一个多用表欧姆档内部电路示意图。电流表满偏电流 $0.5\UmA$、内阻 $10\UO$；电池电动势 $1.5\UV$、内阻 $1\UO$；变阻器 $R_{0}$ 阻值 $0\sim5000\UO$。
\begin{enumerate}
	\item
	该欧姆表的刻度值是按电池电动势为 $1.5\UV$ 刻度的，当电池的电动势下降到 $1.45\UV$、内阻增大到 $4\UO$ 时仍可调零。调零后 $R_{0}$ 阻值将变\tkanswer[1.2cm]{小}（选填“大”或“小”）；若测得某电阻阻值为 $300\UO$，则这个电阻的真实值是\tkanswer[1.8cm]{$290\UO$}。

	\item
	若该欧姆表换了一个电动势为 $1.5\UV$、内阻为 $10\UO$ 的电池，调零后测量某电阻的阻值，其测量结果\tkanswer[1.4cm]{准确}（选填“偏大”、“偏小”或“准确”）。
\end{enumerate}

\onepicture[h=4.2cm, align=right]{\includesvg{figs/27}}
%% answer: （1）小，$290\UO$；（2）准确
\jdanswer{
\begin{enumerate}
	\item
	小，$290\UO$

	\item
	准确

\end{enumerate}
}
%% memo:
\memoanswer{
（1）由欧姆表的调零原理知 $\frac{E}{0.5\UmA}=R_{0}+10\UO+r$，可知 $E=1.5\UV$、$r=1\UO$ 时 $R_{0}=2989\UO$，$E=1.45\UV$、$r=4\UO$ 时，$R_{0}=2886\UO$，调零后 $R_{0}$ 阻值将变小；设测得 $R_{x}=300\UO$ 时对应的电流表偏转为 $I_{x}$，则 $1.5\UV=I_{x}(3000\UO+300\UO)$，$1.45\UV=I_{x}(2900\UO+R_{\text{真}})$，可得 $R_{x}$ 的真实值为 $R_{\text{真}}=290\UO$。

（2）由（1）的计算可知欧姆表内的电池电动势变化影响测量的准确性，电池内阻变化不影响测量的准确性，所以测量结果准确。
}

\item
%% number: 28
%% paperName: 2015年上海高考第 28 题 8 分
%% typeId: 4
%% type: 实验题
%% chapter: 力物体的平衡
%% point: 力矩盘
%% method: 无
%% score: 8
%% degree: 450
%% duplicateId: 0
%% body:
改进后的“研究有固定转动轴物体平衡条件”的实验装置如图所示，力传感器、定滑轮固定在横杆上，替代原装置中的弹簧秤。已知力矩盘上各同心圆的间距均为 $5\Ucm$。
\begin{enumerate}
	\item
	（多选题）做这样改进的优点是

	\fourchoices[answer=BD]
	{力传感器既可测拉力又可测压力}
	{力传感器测力时不受主观判断影响，精度较高}
	{能消除转轴摩擦引起的实验误差}
	{保证力传感器所受拉力方向不变}

	\item
	某同学用该装置做实验，检验时发现盘停止转动时 G 点始终在最低处，他仍用该盘做实验。在对力传感器进行调零后，用力传感器将力矩盘的 G 点拉到图示位置，此时力传感器读数为 $3\UN$。再对力传感器进行调零，然后悬挂钩码进行实验。此方法\tkanswer[1.2cm]{能}（选填“能”、“不能”）消除力矩盘偏心引起的实验误差。已知每个钩码所受重力为 $1\UN$，力矩盘按图示方式悬挂钩码后，力矩盘所受顺时针方向的合力矩为\tkanswer[1.8cm]{$0.7\UNcdotm$}，力传感器的读数为\tkanswer[1.8cm]{$-0.5\UN$}。
\end{enumerate}

\onepicture[h=5.4cm, align=right]{\includesvg{figs/28}}
%% answer: （1）BD；（2）能，$0.7\UNcdotm$，$-0.5\UN$
\jdanswer{
\begin{enumerate}
	\item
	BD

	\item
	能，$0.7\UNcdotm$，$-0.5\UN$

\end{enumerate}
}
%% memo:
\memoanswer{
（1）力传感器只能测量拉力，故 A 错误；力传感器不用实验者目测读数，不受主观判断影响，精度较高，故 B 正确；该装置无法消除转轴摩擦引起的实验误差，故 C 错误；力传感器、定滑轮固定在横杆上，保证力传感器所受拉力方向不变，故 D 正确。

（2）能消除力矩盘偏心带来的误差；顺时针合力矩为 $2\times0.05\UNcdotm+2\times0.15\UNcdotm+3\times0.1\UNcdotm=0.7\UNcdotm$；此时绳上拉力为 $\frac{0.7-3\times0.15}{0.1}\UN=2.5\UN$，故传感器读数为 $2.5\UN-3\UN=-0.5\UN$。
}

\item
%% number: 29
%% paperName: 2015年上海高考第 29 题 9 分
%% typeId: 4
%% type: 实验题
%% chapter: 气体性质
%% point: 盖吕萨克定律
%% method: 无
%% score: 9
%% degree: 250
%% duplicateId: 0
%% body:
简易温度计构造如图所示。两内径均匀的竖直玻璃管下端与软管连接，在管中灌入液体后，将左管上端通过橡皮塞插入玻璃泡。在标准大气压下，调节右管的高度，使左右两管的液面相平，在左管液面位置标上相应的温度刻度。多次改变温度，重复上述操作。
\begin{enumerate}
	\item
	（单选题）此温度计的特点是

	\fourchoices[answer=A]
	{刻度均匀，刻度值上小下大}
	{刻度均匀，刻度值上大下小}
	{刻度不均匀，刻度值上小下大}
	{刻度不均匀，刻度值上大下小}

	\item
	（多选题）影响这个温度计灵敏度的因素有

	\fourchoices[answer=BC]
	{液体密度}
	{玻璃泡大小}
	{左管内径粗细}
	{右管内径粗细}

	\item
	若管中液体是水银，当大气压变为 $75\UcmHg$ 时，用该温度计测得的温度值\tkanswer[1.2cm]{偏大}（选填“偏大”或“偏小”）。为测得准确的温度，在测量时需\tkanswer[4.6cm]{调整两管液面高度差，使右管液面比左管液面高出 $1\Ucm$，然后读数}。
\end{enumerate}

\onepicture[h=5.0cm, align=right]{\includesvg{figs/29}}
%% answer: （1）A；（2）BC；（3）偏大，调整两管液面高度差，使右管液面比左管液面高出 $1\Ucm$，然后读数
\jdanswer{
\begin{enumerate}
	\item
	A

	\item
	BC

	\item
	偏大，调整两管液面高度差，使右管液面比左管液面高出 $1\Ucm$，然后读数

\end{enumerate}
}
%% memo:
\memoanswer{
（1）实验中调节左右两管液面相平，则玻璃泡内压强不变，始终等于大气压强，由气体状态方程可知，左管内密闭气体的体积与温度成正比，温度越高，体积越大，左管液面越低，所以刻度值为上小下大，故 A 正确。

（2）温度改变时，左管高度改变得越大，灵敏度越高，即灵敏度可以表示为 $\frac{\Delta h}{\Delta T}$。由等压变化规律可知 $\frac{\Delta V}{\Delta T}=\frac{\Delta h}{\Delta T}S=\frac{V}{T}$，得 $\frac{\Delta h}{\Delta T}=\frac{V}{ST}$，式中 $S$ 为左管的横截面积。所以左管内径越细，玻璃泡体积 $V$ 越大灵敏度越高。右管与大气连通，其粗细对灵敏度没有影响，故 BC 正确。

（3）大气压强变小，左管液面会降低，测量的温度偏大；为了测量准确，可使右管比左管液面恰好高 $1\Ucm$ 时读数，保证左管内密闭气体的压强等于标准大气压。
}

\item
%% number: 30
%% paperName: 2015年上海高考第 30 题 10 分
%% typeId: 6
%% type: 计算题
%% chapter: 气体性质
%% point: 气体连接体
%% method: 无
%% score: 10
%% degree: 650
%% duplicateId: 0
%% body:
如图，气缸左右两侧气体由绝热活塞隔开，活塞与气缸光滑接触。初始时两侧气体均处于平衡态，体积之比 $V_{1}:V_{2}=1:2$，温度之比 $T_{1}:T_{2}=2:5$。先保持右侧气体温度不变，升高左侧气体温度，使两侧气体体积相同；然后使活塞导热，两侧气体最后达到平衡。求：
\begin{enumerate}
	\item
	两侧气体体积相同时，左侧气体的温度与初始温度之比；

	\item
	最后两侧气体的体积之比。
\end{enumerate}

\onepicture[h=2.6cm, align=right]{\includesvg{figs/30}}
%% answer: （1）2；（2）$5:4$
\jdanswer{
\begin{enumerate}
	\item
	2

	\item
	$5:4$

\end{enumerate}
}
%% memo:
\memoanswer{
（1）设初始时压强为 $p$。左侧气体满足 $\frac{pV_{1}}{T_{1}}=\frac{p'V}{kT_{1}}$，右侧气体满足 $pV_{2}=p'V$，解得 $k=\frac{V_{2}}{V_{1}}=2$。

（2）活塞导热达到平衡，左侧气体满足 $\frac{p'V}{kT_{1}}=\frac{p''V_{1}'}{T_{1}'}$，右侧气体满足 $\frac{p'V}{T_{2}}=\frac{p''V_{2}'}{T_{2}'}$，平衡时 $T_{1}'=T_{2}'$，解得 $\frac{V_{1}'}{V_{2}'}=\frac{T_{2}}{kT_{1}}=\frac{5}{4}$。
}

\item
%% number: 31
%% paperName: 2015年上海高考第 31 题 12 分
%% typeId: 6
%% type: 计算题
%% chapter: 运动和力
%% point: 加速减速模型
%% method: 无
%% score: 12
%% degree: 450
%% duplicateId: 0
%% body:
质量为 $m$ 的小球在竖直向上的恒定拉力作用下，由静止开始从水平地面向上运动，经一段时间，拉力做功为 $W$，此后撤去拉力，球又经相同时间回到地面。以地面为零势能面，不计空气阻力。求：
\begin{enumerate}
	\item
	球回到地面时的动能 $E_{kt}$；

	\item
	撤去拉力前球的加速度大小 $a$ 及拉力的大小 $F$；

	\item
	球动能为 $\frac{W}{5}$ 时的重力势能 $E_{p}$。
\end{enumerate}
%% answer: （1）$W$；（2）$a=\frac{1}{3}g$，$F=\frac{4}{3}mg$；（3）$\frac{3}{5}W$ 或 $\frac{4}{5}W$
\jdanswer{
\begin{enumerate}
	\item
	$W$

	\item
	$a=\frac{1}{3}g$，$F=\frac{4}{3}mg$

	\item
	$\frac{3}{5}W$ 或 $\frac{4}{5}W$

\end{enumerate}
}
%% memo:
\memoanswer{
（1）撤去拉力时球的机械能为 $W$，由机械能守恒，回到地面时的动能 $E_{kt}=W$。

（2）设拉力作用时间为 $t$，在此过程中球上升 $h$，末速度为 $v$，则 $h=\frac{1}{2}at^{2}$，$v=at$，由题意有 $-h=vt-\frac{1}{2}gt^{2}$，解得 $a=\frac{1}{3}g$，根据牛顿定律 $F-mg=ma$，解得 $F=\frac{4}{3}mg$。

（3）动能为 $\frac{W}{5}$ 时球的位置可能在 $h$ 的下方或上方。设球的位置在 $h$ 下方离地 $h'$ 处，$(F-mg)h'=\frac{1}{5}W$，而 $(F-mg)h=\frac{1}{4}W$，解得 $h'=\frac{4}{5}h$，重力势能 $E_{p}=mgh'=\frac{3}{5}W$。设球的位置在 $h$ 上方离地 $h''$ 处，由机械能守恒 $\frac{1}{5}W+mgh''=W$，因此 $E_{p}=mgh''=\frac{4}{5}W$。
}

\item
%% number: 32
%% paperName: 2015年上海高考第 32 题 14 分
%% typeId: 6
%% type: 计算题
%% chapter: 电磁感应
%% point: 电磁感应受力分析
%% method: 无
%% score: 14
%% degree: 450
%% duplicateId: 0
%% body:
如图（a），两相距 $L=0.5\Um$ 的平行金属导轨固定于水平面上，导轨左端与阻值 $R=2\UO$ 的电阻连接，导轨间虚线右侧存在垂直导轨平面的匀强磁场。质量 $m=0.2\Ukg$ 的金属杆垂直置于导轨上，与导轨接触良好，导轨与金属杆的电阻可忽略。杆在水平向右的恒定拉力作用下由静止开始运动，并始终与导轨垂直，其 $v-t$ 图像如图（b）所示。在 $15\Us$ 时撤去拉力，同时使磁场随时间变化，从而保持杆中电流为 0。求：
\begin{enumerate}
	\item
	金属杆所受拉力的大小 $F$；

	\item
	$0\sim15\Us$ 内匀强磁场的磁感应强度大小 $B_{0}$；

	\item
	$15\sim20\Us$ 内磁感应强度随时间的变化规律。
\end{enumerate}

\twopicture[numA={（a）}, numB={（b）}, labelA=2015上海32a, labelB=2015上海32b, hA=3.0cm, hB=3.0cm]
{\includesvg[inkscapearea=page]{figs/32a}}
{\includesvg[inkscapearea=page]{figs/32b}}

%% answer: （1）$F=0.24\UN$；（2）$B_{0}=0.4\UT$；（3）$B_{t}=\frac{20}{50-(t-15)(t-25)}\UT$
\jdanswer{
\begin{enumerate}
	\item
	$F=0.24\UN$

	\item
	$B_{0}=0.4\UT$

	\item
	$B_{t}=\frac{20}{50-(t-15)(t-25)}\UT$

\end{enumerate}
}
%% memo:
\memoanswer{
（1）由 $v-t$ 关系图可知在 $0\sim10\Us$ 时间段杆尚未进入磁场，因此 $F-\mu mg=ma_{1}$，由图可得 $a_{1}=0.4\Umsq$，同理可知在 $15\sim20\Us$ 时间段杆仅在摩擦力作用下运动。$\mu mg=ma_{2}$，由图可得 $a_{2}=0.8\Umsq$，解得 $F=0.24\UN$。

（2）在 $10\sim15\Us$ 时间段杆在磁场中做匀速运动，因此有 $F=\mu mg+\frac{B_{0}^{2}L^{2}v}{R}$，以 $F=0.24\UN$，$\mu mg=0.16\UN$ 代入，解得 $B_{0}=0.4\UT$。

（3）由题意可知在 $15\sim20\Us$ 时间段通过回路的磁通量不变，设杆在 $10\sim15\Us$ 内运动距离为 $d$，$15\Us$ 后运动距离为 $x$，$B_{t}L(d+x)=B_{0}Ld$，其中 $d=20\Um$，$x=4(t-15)-0.4(t-15)^{2}$，由此可得 $B_{t}=\frac{B_{0}d}{d+x}=\frac{20}{50-(t-15)(t-25)}\UT$。
}

\item
%% number: 33
%% paperName: 2015年上海高考第 33 题 14 分
%% typeId: 6
%% type: 计算题
%% chapter: 电场
%% point: 带电粒子的圆周运动
%% method: 分情况讨论
%% score: 14
%% degree: 250
%% duplicateId: 0
%% body:
如图，在场强大小为 $E$、水平向右的匀强电场中，一轻杆可绕固定转轴 O 竖直平面内自由转动。杆的两端分别固定两电荷量均为 $q$ 的小球 A、B，A 带正电，B 带负电；A、B 两球到转轴 O 的距离分别为 $2l$、$l$，所受重力大小均为电场力大小的 $\sqrt{3}$ 倍。开始时杆与电场间夹角为 $\theta$（$90^{\circ}\leq\theta\leq180^{\circ}$）。将杆从初始位置由静止释放，以 O 点为重力势能和电势能零点。求：
\begin{enumerate}
	\item
	初始状态的电势能 $W_{e}$；

	\item
	杆在平衡位置时与电场间的夹角 $\alpha$；

	\item
	杆在电势能为零处的角速度 $\omega$。
\end{enumerate}

\onepicture[h=2.6cm, align=right]{\includesvg{figs/33}}
%% answer: （1）$W_{e}=-3Eql\cos\theta$；（2）$\alpha=30^{\circ}$；（3）当 $\theta<150^{\circ}$ 时，$\omega_{1}=\sqrt{\frac{2g}{5l}(1-\sin\theta-\sqrt{3}\cos\theta)}$；当 $\theta\geq150^{\circ}$ 时有两个位置，$\omega_{1}=\sqrt{\frac{2g}{5l}(1-\sin\theta-\sqrt{3}\cos\theta)}$，$\omega_{2}=\sqrt{\frac{2g}{5l}(-1-\sin\theta-\sqrt{3}\cos\theta)}$
\jdanswer{
\begin{enumerate}
	\item
	$W_{e}=-3Eql\cos\theta$

	\item
	$\alpha=30^{\circ}$

	\item
	当 $\theta<150^{\circ}$ 时，$\omega_{1}=\sqrt{\frac{2g}{5l}(1-\sin\theta-\sqrt{3}\cos\theta)}$；当 $\theta\geq150^{\circ}$ 时有两个位置，$\omega_{1}=\sqrt{\frac{2g}{5l}(1-\sin\theta-\sqrt{3}\cos\theta)}$，$\omega_{2}=\sqrt{\frac{2g}{5l}(-1-\sin\theta-\sqrt{3}\cos\theta)}$

\end{enumerate}
}
%% memo:
\memoanswer{
（1）初态时 $W_{e}=q\varphi_{A}+(-q)\varphi_{B}=q(\varphi_{A}-\varphi_{B})=qU_{AB}=-3Eql\cos\theta$。

（2）平衡位置如图。设小球质量为 $m$，合力矩为 $3Eql\sin\alpha-mgl\cos\alpha=0$，由此得 $\tan\alpha=\frac{mg}{3qE}=\frac{\sqrt{3}}{3}$，$\alpha=30^{\circ}$。

\onepicture[h=3.0cm]{\includesvg{figs/33_an}}

（3）由分析知，杆在电势能为零处时杆与电场方向垂直，可能 A 球在上、B 球在下，也可能 B 球在上、A 球在下。设杆与电场间夹角为 $\theta_{0}$ 时，A 球恰能到达 O 正上方，在此位置杆的角速度为 0。初态：$W_{e}=-3Eql\cos\theta_{0}$，$E_{p}=-mgl\sin\theta_{0}$；末态：$W_{e}'=0$，$E_{p}'=mgl$。根据能量守恒有：$-3Eql\cos\theta_{0}-mgl\sin\theta_{0}=mgl$，可得 $\theta_{0}=150^{\circ}$。当 $\theta<150^{\circ}$ 时，A 球无法到达最高点。当 A 位于 O 正下方处电势能为零。初态：$W_{e}=-3Eql\cos\theta$，$E_{p}=-mgl\sin\theta$，$E_{k}=0$；末态：$W_{e}'=0$，$E_{p}'=-mgl$，$E_{k}'=\frac{1}{2}m(\omega l)^{2}+\frac{1}{2}m(\omega\cdot2l)^{2}=\frac{5}{2}ml^{2}\omega^{2}$。根据能量守恒有：$-3Eql\cos\theta-mgl\sin\theta=\frac{5}{2}ml^{2}\omega^{2}-mgl$，$\omega=\sqrt{\frac{2mg(1-\sin\theta)-6qE\cos\theta}{5ml}}=\sqrt{\frac{2g}{5l}(1-\sin\theta-\sqrt{3}\cos\theta)}$。当 $\theta\geq150^{\circ}$ 时，电势能为零的位置有两处，即 A 位于 O 正下方或正上方处。在 A 位于 O 正下方时 $\omega=\sqrt{\frac{2g}{5l}(1-\sin\theta-\sqrt{3}\cos\theta)}$。在 A 位于 O 正上方时，有 $-3Eql\cos\theta-mgl\sin\theta=\frac{5}{2}ml^{2}\omega^{2}+mgl$，$\omega=\sqrt{\frac{-2mg(1+\sin\theta)-6qE\cos\theta}{5ml}}=\sqrt{\frac{2g}{5l}(-1-\sin\theta-\sqrt{3}\cos\theta)}$。
}

\end{enumerate}

\end{document}
